\begin{document}

\showtodostrue
\showcommentstrue

\newauthorcomment{alice}{Alice}{red}
\newauthorcomment{bob}{Bob}{red}

\subtitle{Lecture notes}

% Optional cover figure: remove this command for a plain cover.
\titlefigure{%
    \begin{tikzpicture}[
        x={(1cm,0cm)},
        y={(0.45cm,0.85cm)}
    ]
        \filldraw[
            fill=CustomBlue!8,
            draw=CustomBlue,
            line width=0.7pt
        ]
            (0,0) -- (1,0) -- (1,1) -- (0,1) -- cycle;

        \foreach \i in {-2,...,3} {
            \foreach \j in {-1,...,2} {
                \fill[black!65] (\i,\j) circle[radius=1.2pt];
            }
        }
    \end{tikzpicture}%
}

\hypersetup{pageanchor=false}
\maketitle

% Document information page
\thispagestyle{empty}
\noindent
These are lecture notes for an undergraduate course on lattice theory.

\vspace{2\baselineskip}

\noindent\textbf{Course information}\par
\smallskip
\noindent
Undergraduate course on lattice theory, Example University.\\
Academic year 2026--2027.

\vspace{\baselineskip}

\noindent\textbf{Status of these notes}\par
\smallskip
\noindent
This is a working document. Corrections and suggestions are welcome.

\par
\vfill
\printauthorinfo
\clearpage

\pagenumbering{roman}
\hypersetup{pageanchor=true}
\tableofcontents
\clearpage

\pagenumbering{arabic}

\chapter{A Sample Chapter}

\section{Structured statements}

\begin{definition}
    \label{definition:longform-object}
    A \emph{longform object} is an object \(X\) equipped with a map
    \[
        \alpha_X \colon X \to G(X)
    \]
    satisfying the following properties.
    \begin{theoremitems}
    \item The object \(X\) exists in the category under consideration.
        \label{item:longform-existence}
    \item It is unique up to unique isomorphism.
        \label{item:longform-uniqueness} 
    \end{theoremitems}
\end{definition}

\begin{proposition}
    \label{proposition:longform-main}
    Every longform object is determined by the data in
    \Cref{item:longform-existence,item:longform-uniqueness}.
\end{proposition}

\begin{proof}
    This is a template proof. The point is that the item references are attached to
    the theorem-like environment \Cref{definition:longform-object}, so their labels
    are numbered using the surrounding theorem number.
\end{proof}

\begin{remark}
    \label{remark:longform-citations}
    Long-form documents often contain many references.
    This example cites Hartshorne \cite{har77} and the Stacks Project \cite{stacks}.
\end{remark}

\section{Diagrams and ordinary lists}

The template supports commutative diagrams. For instance, one may write
\[
    \begin{tikzcd}
        A \arrow[r, "\phi"] \arrow[d, "\alpha"'] &
        B \arrow[d, "\beta"] \\
        C \arrow[r, "\psi"'] &
        D .
    \end{tikzcd}
\]
Such diagrams are useful in lecture notes and book-length projects.

The template also supports ordinary enumerated lists:
\begin{enumerate}[label=(\alph*), ref=(\alph*)]
    \item A first ordinary item.
        \label{item:longform-ordinary-first} 
    \item A second ordinary item.
        \label{item:longform-ordinary-second} 
\end{enumerate}
We can refer to \Cref{item:longform-ordinary-second} as an ordinary item, and
to \Cref{item:longform-uniqueness} as a theorem-numbered item.

\chapter{Another Chapter}

\section{Resetting theorem numbers}

This chapter checks that theorem-like numbering behaves correctly after a
chapter change.
\todo{Here is a to-do note.}

\begin{lemma}
    \label{lemma:second-chapter}
    The references from the first chapter, such as
    \Cref{definition:longform-object}, still work from later chapters.
\end{lemma}

\begin{proof}
    This proof is only illustrative. It also demonstrates a displayed equation:
    \begin{equation}
        \label{equation:pull-back}
        (f \circ g)^* = g^* \circ f^* .
    \end{equation}
\end{proof}

\alice{We should add an example here.}

\begin{example}
    \label{example:external-link}
    A long-form document may also contain external links, for instance to
    \href{https://www.ctan.org/pkg/biblatex}{the biblatex package page}.
\end{example}

\bob{I've added the example above.}

\printbibliography

\end{document}
